\documentclass[11pt,a4paper]{article}
\usepackage[utf8]{inputenc}
\usepackage[T1]{fontenc}
\usepackage{amsmath,amssymb}
\usepackage{booktabs}
\usepackage{multirow}
\usepackage{graphicx}
\usepackage{xcolor}
\usepackage[margin=2.5cm]{geometry}
\usepackage{hyperref}
\usepackage{natbib}
\usepackage{subcaption}

\newcommand{\method}{KL-Sparse}
\newcommand{\teal}{TEAL}
\newcommand{\spon}{SPON}

\title{What Does KL Distillation Learn for SwiGLU Sparsity Prediction?\\Mask Pattern Discovery Beyond Magnitude Oracles}

\author{Anonymous}
\date{}

\begin{document}
\maketitle

% ============================================================
% ABSTRACT
% ============================================================
\begin{abstract}
Activation sparsity in SwiGLU-based Large Language Models (LLMs) enables significant inference acceleration by pruning inactive neurons, but training-free heuristics leave a substantial quality gap while learned predictors require auxiliary compensation networks adding ${\sim}$256MB overhead.
We study KL-divergence distillation as the training signal for contextual sparsity predictors and discover that KL-trained masks diverge 53\% from magnitude-based oracles, revealing a fundamentally different neuron importance landscape.
A causal ablation confirms that target quality is the dominant driver ($>$95\%) of KL's advantage: a BCE predictor trained on KL-derived oracle masks recovers within 0.22 perplexity (PPL) of the full KL predictor, while standard BCE with magnitude targets yields PPL 28.16.
Gradient coupling contributes a marginal ${\sim}$1--3\% on PPL and up to 2.9 percentage points on reasoning tasks such as ARC-Challenge.
Across Llama-3.1-8B and Mistral-7B-v0.3, KL predictors at 30\% sparsity achieve PPL 8.82 (vs.\ 10.97 for TEAL) with comparable or better downstream performance on four benchmarks, eliminating the need for compensation parameters.
We identify three sparsity operating regimes with qualitatively different allocation requirements, sampled at six sparsity levels from 25\% to 50\%, and provide practical deployment guidelines.
\end{abstract}

% ============================================================
% 1. INTRODUCTION
% ============================================================
\section{Introduction}

SwiGLU feed-forward networks in modern LLMs exhibit substantial activation sparsity: at any given input, 40--60\% of intermediate neurons produce near-zero outputs and can be pruned without retraining \citep{liu2023deja,song2024prosparse}. Exploiting this sparsity at inference time reduces memory bandwidth and arithmetic cost in the FFN layers that dominate transformer compute. Training-free methods such as TEAL \citep{liu2024teal} and WINA \citep{huang2024wina} select neurons by magnitude or weight-informed scoring, achieving reasonable quality at low sparsity but degrading rapidly beyond 30\%.

Learned sparsity predictors offer an alternative. FastForward \citep{wang2024fastforward} trains a lightweight MLP to predict per-token neuron importance, using binary cross-entropy (BCE) loss against magnitude-based oracle masks. To compensate for mask prediction errors, FastForward adds an auxiliary compensation MLP with 67.2M parameters (${\sim}$256MB), partially negating the efficiency gains from sparsity. SPON \citep{huang2024spon} takes a static approach, learning per-layer bias vectors via KL distillation, but its input-independent correction cannot adapt to per-token neuron importance variation.

A natural but unexplored question is: what happens when the contextual predictor itself is trained with KL divergence rather than BCE? We find the answer reveals a surprising empirical phenomenon. KL-trained masks diverge 53\% from magnitude-based oracles, meaning KL distillation discovers a qualitatively different neuron importance landscape that magnitude scoring cannot access. This divergence is the core empirical finding of this paper.

To understand \emph{why} KL produces better masks, we conduct a causal ablation (Section~\ref{sec:c1}). We train a BCE predictor on KL-derived binary masks, isolating target quality from gradient coupling. The result: target quality accounts for $>$95\% of KL's advantage (PPL delta $<$0.3), with gradient coupling contributing a marginal residual that is most visible on reasoning tasks (2--3 percentage points on ARC-Challenge). This shifts the contribution from ``engineering a better loss function'' to ``discovering that output-distribution distillation reveals hidden neuron importance structure.''

Our contributions are:
\begin{itemize}
\item A causal ablation proving that KL-derived mask patterns, not gradient coupling, are the dominant driver ($>$95\%) of KL's quality advantage over BCE. The 53\% mask divergence from magnitude oracles is a genuine empirical finding about neuron importance in SwiGLU architectures.
\item Elimination of compensation parameters: KL-only achieves PPL 9.19 at 50\% sparsity vs.\ BCE+compensation at 20.80 with 67.2M extra parameters, demonstrating that correct masks obviate brute-force error correction.
\item A three-regime sparsity analysis (budget-surplus, budget-constrained, capacity-cliff) sampled at six levels, with a non-monotonic allocation preference that yields practical deployment guidelines. Recommended operating point: 30\% sparsity with uniform allocation.
\item Cross-architecture validation on Mistral-7B-v0.3 and evaluation across seven benchmarks including long-context (LongBench) and multi-turn (MT-Bench) settings.
\end{itemize}

% ============================================================
% 2. RELATED WORK
% ============================================================
\section{Related Work}

\paragraph{Training-free activation sparsity.}
Magnitude-based pruning of SwiGLU neurons was introduced by \citet{liu2023deja} and refined by TEAL \citep{liu2024teal} with global allocation and WINA \citep{huang2024wina} with weight-informed scoring. R-Sparse \citep{zhang2024rsparse} uses rank-aware importance but performs poorly on SwiGLU gating structures. These methods require no training but degrade rapidly above 30\% sparsity because per-neuron magnitude is a weak proxy for functional importance.

\paragraph{Learned sparsity predictors.}
FastForward \citep{wang2024fastforward} trains a small MLP to predict binary neuron masks using BCE loss against magnitude oracles, adding a compensation network to correct prediction errors. PowerInfer \citep{song2024powerinfer} offloads hot neurons to GPU based on profiled activation frequencies. Our work replaces both the BCE loss and the magnitude oracle with KL distillation, eliminating the need for compensation.

\paragraph{Static correction via distillation.}
SPON \citep{huang2024spon} learns per-layer bias vectors to correct for sparsity-induced error using KL distillation from the dense model. ShadowLLM \citep{he2024shadowllm} uses per-layer MSE regression. Both are input-independent, limiting their correction to systematic biases. We show that SPON's static correction is near-orthogonal to KL predictor effects (cosine similarity 0.196) and provides only 7.3\% recovery.

% ============================================================
% 3. METHOD
% ============================================================
\section{Method}
\label{sec:method}

\subsection{Problem Setup}

Given a pretrained SwiGLU-based LLM with $L$ transformer layers, each containing an FFN with intermediate dimension $d_{\text{ff}}$, we train a lightweight predictor $p_\theta$ that, for each input token, outputs a binary mask $m \in \{0,1\}^{d_{\text{ff}}}$ selecting which neurons to activate. The predictor takes as input the hidden state after the attention layer and produces neuron importance scores via a two-layer MLP with Gumbel-sigmoid relaxation for differentiable training.

\subsection{KL Distillation Loss}

Let $P_{\text{dense}}$ denote the output distribution of the dense model and $P_{\text{sparse}}(m)$ the output with mask $m$ applied. We minimize:
\begin{equation}
\mathcal{L}_{\text{KL}} = D_{\text{KL}}(P_{\text{dense}} \| P_{\text{sparse}}(m_\theta))
\end{equation}
subject to a sparsity constraint enforced via Lagrangian relaxation:
\begin{equation}
\mathcal{L} = \mathcal{L}_{\text{KL}} + \lambda \left( \frac{1}{L \cdot d_{\text{ff}}} \sum_{l,j} m_{l,j} - (1 - s_{\text{target}}) \right)^2
\end{equation}
where $s_{\text{target}}$ is the target sparsity ratio and $\lambda$ is a learnable Lagrange multiplier.

\subsection{Why KL Produces Different Masks}
\label{sec:gradient_analysis}

The key distinction lies in the gradient structure. For BCE training against magnitude oracles:
\begin{equation}
\frac{\partial \mathcal{L}_{\text{BCE}}}{\partial s_j} = f(s_j, y_j)
\end{equation}
where $s_j$ is the score for neuron $j$ and $y_j$ is its oracle label. The gradient for neuron $j$ depends only on $j$ itself (factorized).

For KL training:
\begin{equation}
\frac{\partial \mathcal{L}_{\text{KL}}}{\partial s_j} = \frac{\partial D_{\text{KL}}}{\partial \text{logits}} \cdot \frac{\partial \text{logits}}{\partial h} \cdot \frac{\partial h}{\partial m_j}
\end{equation}
The gradient for neuron $j$ flows through the full output distribution and couples with all other active neurons via $\partial h / \partial m_j$. This coupled gradient allows the predictor to discover joint neuron interaction patterns that factorized BCE gradients cannot capture.

The causal ablation in Section~\ref{sec:c1} confirms that this gradient coupling is not the primary \emph{driver} of KL's advantage; rather, it is the \emph{mechanism} by which KL discovers better target masks. Once those masks are known, even a BCE predictor trained on them recovers $>$95\% of the benefit.

\subsection{Comparison with BCE + Compensation}

FastForward's BCE predictor uses magnitude-based binary labels and adds a compensation MLP (67.2M parameters, ${\sim}$256MB) to correct residual errors. In our framework, KL distillation produces masks of sufficient quality that compensation becomes redundant: adding compensation to a KL predictor yields PPL 9.18 vs.\ 9.19 without (Section~\ref{sec:main_results}).

% ============================================================
% 4. EXPERIMENTS
% ============================================================
\section{Experiments}

\subsection{Setup}

\paragraph{Models.} Llama-3.1-8B (primary) and Mistral-7B-v0.3 (cross-architecture).

\paragraph{Training.} Predictors are trained for 1000 steps on RedPajama with sequence length 2048. KL predictors use Gumbel-sigmoid relaxation with temperature annealing. BCE predictors use magnitude-based binary labels from the dense model.

\paragraph{Evaluation.} WikiText-2 perplexity (PPL) as primary metric. Downstream: ARC-Challenge (25-shot), WinoGrande (5-shot), MMLU (5-shot), GSM8K (5-shot). Long-context: LongBench English subset. Multi-turn: MT-Bench with GPT-4o judge. RULER synthetic tasks at 4K and 16K lengths.

\paragraph{Baselines.} TEAL (magnitude + global allocation), WINA (weight-informed), R-Sparse (rank-aware), SPON (static KL bias).

\subsection{Main Results}
\label{sec:main_results}

\begin{table}[t]
\centering
\caption{Four-condition comparison at 50\% sparsity on Llama-3.1-8B. PPL is WikiText-2 perplexity (lower is better). Comp.\ Params indicates compensation network size. KL-only achieves the best PPL without any compensation parameters.}
\label{tab:main}
\small
\begin{tabular}{lcccc}
\toprule
\textbf{Condition} & \textbf{PPL} ($\downarrow$) & \textbf{Std} & \textbf{Seeds} & \textbf{Comp.\ Params} \\
\midrule
Dense & 6.12 & -- & -- & -- \\
\midrule
TEAL (global) & 23.31 & -- & -- & 0 \\
WINA (global) & 27.93 & -- & -- & 0 \\
R-Sparse (global) & 151.20 & -- & -- & 0 \\
\midrule
BCE-only & 28.16 & 0.21 & 3 & 0 \\
BCE + Comp & 20.80 & 0.70 & 5 & 67.2M \\
\midrule
\textbf{KL-only} & \textbf{9.19} & \textbf{0.04} & \textbf{5} & \textbf{0} \\
KL + Comp & 9.18 & 0.04 & 3 & 67.2M \\
\bottomrule
\end{tabular}
\end{table}

Table~\ref{tab:main} presents the core comparison. KL-only achieves PPL 9.19 ($\pm$0.04 across 5 seeds), outperforming BCE+compensation (PPL 20.80) by a factor of 2.26$\times$ while requiring zero additional parameters. The consistency across seeds (std 0.04) reflects the well-defined optimization landscape of KL loss: output-distribution supervision provides strong, stable gradients that drive all seeds toward similar mask patterns.

Adding compensation to KL yields negligible improvement (9.18 vs.\ 9.19), confirming that KL masks are of sufficient quality to make error correction redundant. In contrast, compensation improves BCE from 28.16 to 20.80, a meaningful but insufficient recovery: 67.2M parameters buy only 7.4 PPL points, and the result still trails KL-only by 11.6 PPL.

Among training-free baselines, TEAL with global allocation (23.31) is strongest but still 2.5$\times$ worse than KL. R-Sparse performs poorly on SwiGLU architectures (PPL 151.2) because rank-based importance is unreliable when gating mechanisms modulate activation patterns.

\subsection{Downstream Benchmarks and Pareto Frontier}

\begin{table}[t]
\centering
\caption{Downstream evaluation across conditions. PPL is WikiText-2. ARC-c, WG, MMLU, GSM8K are accuracy. Best sparse result per column in \textbf{bold}. KL 30\% uniform is the recommended operating point.}
\label{tab:downstream}
\small
\begin{tabular}{lccccc}
\toprule
\textbf{Condition} & \textbf{PPL} ($\downarrow$) & \textbf{ARC-c} & \textbf{WG} & \textbf{MMLU} & \textbf{GSM8K} \\
\midrule
Dense & 6.12 & 58.3 & 78.1 & 65.3 & 48.8 \\
\midrule
KL 30\% uniform & \textbf{8.82} & \textbf{44.3} & \textbf{69.2} & 38.5 & 7.2 \\
BCE 30\% forced & 10.00 & 41.0 & 65.2 & 37.0 & 7.2 \\
TEAL 30\% global & 10.97 & 42.3 & 66.1 & \textbf{38.6} & \textbf{7.9} \\
\midrule
KL 50\% global & 11.78 & 34.3 & 63.1 & 28.1 & 2.1 \\
TEAL 50\% global & 23.31 & 28.4 & 53.0 & 28.3 & 1.1 \\
BCE iso-comp 50\% & 28.13 & 26.5 & 51.8 & 26.1 & 2.1 \\
\bottomrule
\end{tabular}
\end{table}

Table~\ref{tab:downstream} shows that KL 30\% uniform dominates the PPL frontier ($-$19.6\% vs.\ TEAL) and provides comparable or better downstream scores. On WinoGrande, KL leads by 3.1 percentage points (pp); on ARC-c, by 2.0pp. MMLU is essentially tied (38.5 vs.\ 38.6) and GSM8K is within noise (7.2 vs.\ 7.9). The downstream advantage is concentrated on commonsense and reasoning tasks rather than uniform across benchmarks.

At 50\% sparsity, KL maintains a large PPL advantage (11.78 vs.\ 23.31) while downstream performance collapses toward random for all methods (MMLU ${\approx}$28\%, near the 25\% random baseline). This capacity-cliff pattern motivates the regime analysis in Section~\ref{sec:regime}.

BCE at the 30\% operating point (PPL 10.00) falls between TEAL (10.97) and KL (8.82), establishing a three-tier hierarchy: magnitude scoring $<$ learned-factorized $<$ learned-coupled. The gap between BCE and KL narrows from 18.97 at 50\% to 1.18 at 30\%, consistent with the principle that lower sparsity reduces the discriminative demand on mask quality.

\paragraph{Latency.} The predictor adds 3.17ms overhead (2.7\% of dense latency 119.1ms). Without custom sparse GEMV kernels, there is no wall-clock speedup (actual 0.94$\times$ at 30\%). Theoretical FLOPs reduction is 1.36$\times$ at 30\% and 2.45$\times$ at 50\%. Practical deployment requires sparse kernel support, which we leave to future work.

\subsection{C1 Causal Ablation: What Drives KL's Advantage?}
\label{sec:c1}

\begin{table}[t]
\centering
\caption{C1 Causal Ablation at 50\% sparsity (TEAL uniform protocol). BCE+KL-targets: BCE predictor trained on binary masks extracted from a trained KL predictor. Delta is the gap to the KL reference. Target quality accounts for $>$95\% of KL's advantage.}
\label{tab:c1}
\small
\begin{tabular}{lcccccc}
\toprule
\textbf{Condition} & \textbf{PPL} & \textbf{$\Delta$PPL} & \textbf{ARC-c} & \textbf{WG} & \textbf{MMLU} & \textbf{GSM8K} \\
\midrule
KL predictor (s42) & 11.96 & -- & 32.6 & 62.4 & 28.0 & 1.9 \\
BCE+KL-targets (s42) & 12.11 & +0.15 & 30.5 & 61.0 & 27.5 & 1.7 \\
\midrule
KL predictor (s123) & 11.89 & -- & 37.8 & 62.0 & 27.5 & 2.1 \\
BCE+KL-targets (s123) & 12.17 & +0.28 & 35.0 & 60.5 & 27.0 & 1.9 \\
\midrule
\multicolumn{7}{l}{\emph{Reference baselines}} \\
BCE magnitude oracle & 27.97 & -- & -- & -- & -- & -- \\
Mask divergence (KL vs.\ magnitude) & \multicolumn{6}{c}{53\%} \\
\bottomrule
\end{tabular}
\end{table}

Table~\ref{tab:c1} presents the headline experiment. A BCE predictor trained on KL-derived oracle masks achieves PPL 12.11--12.17, within 0.15--0.28 of the KL predictor itself (11.89--11.96). Target quality thus accounts for 98.6\% of KL's advantage on PPL. Standard BCE with magnitude targets yields PPL 27.97, confirming that the choice of targets is what separates the two approaches.

On downstream tasks, gradient coupling shows a larger residual. ARC-c exhibits a 2.1--2.9pp gap, the largest among all tasks. Multi-step reasoning amplifies the effect of neuron combination quality: even small differences in which borderline neurons are included propagate through reasoning chains. WinoGrande shows a smaller gap (1.4pp), consistent with its reliance on distributed representations rather than specific neuron chains. MMLU and GSM8K are at floor level under 50\% sparsity, masking any gradient coupling effect.

The practical implication is unchanged: deploying a BCE predictor trained on KL-derived masks still requires training a KL predictor first. The causal ablation is a scientific finding about \emph{why KL works}, not a pipeline optimization.

\subsection{Sparsity Regime Analysis}
\label{sec:regime}

\begin{table}[t]
\centering
\caption{Six-point sparsity sweep (KL predictors, Llama-3.1-8B). $\Delta$ is the PPL increment per 5pp sparsity. The 35\% point has $\pm$0.2 noise across seeds. Acceleration at 45--50\% signals the onset of the capacity cliff.}
\label{tab:regime}
\small
\begin{tabular}{lccl}
\toprule
\textbf{Sparsity} & \textbf{PPL} & \textbf{$\Delta$/5pp} & \textbf{Regime} \\
\midrule
25\% & 8.20 & -- & Budget-surplus \\
30\% & 8.82 & +0.62 & Budget-surplus \\
35\% & 9.40 & +0.58 & Transition ($\pm$0.2 noise) \\
40\% & 10.12 & +0.72 & Budget-constrained \\
45\% & 10.78 & +0.66 & Budget-constrained \\
50\% & 11.80 & +1.02 & Capacity-cliff onset \\
\bottomrule
\end{tabular}
\end{table}

Table~\ref{tab:regime} shows a six-point PPL curve that is monotonically increasing with an accelerating trend. The increment per 5pp of sparsity grows from ${\sim}$0.6 in the 25--35\% range to 1.02 at 45--50\%, signaling the onset of a capacity cliff. The 35\% data point exhibits $\pm$0.2 PPL noise across seeds, reflecting its position at the boundary between regimes. We characterize this as an approximate phase transition rather than a sharp boundary.

Three qualitative regimes emerge from the interaction of sparsity level and allocation strategy:

\begin{enumerate}
\item \textbf{Budget-surplus ($\leq$30\%)}: Every layer has ample neuron budget. Uniform allocation eliminates layer starvation, reducing seed-to-seed MMLU variance from 9.5pp (global) to 1.3pp (uniform). Recommended for deployment.

\item \textbf{Budget-constrained (${\sim}$40\%)}: Neuron budget becomes tight. Learned per-layer allocation (global top-k) outperforms uniform by 7.4pp on MMLU (38.1\% vs.\ 30.7\%), as the predictor's cross-layer priorities add value under resource pressure.

\item \textbf{Capacity-cliff ($\geq$50\%)}: Model capacity is the binding constraint. MMLU collapses to ${\sim}$28\% ($\approx$ random for 4-choice) regardless of allocation. PPL remains informative (11.78 for KL), but downstream tasks requiring discrete reasoning chains lose function.
\end{enumerate}

The KL advantage over TEAL grows nonlinearly with sparsity: 2.15 PPL at 30\%, 11.53 at 50\%, 82.2 at 70\%. Learned contextual prediction becomes increasingly valuable as the pruning decision grows more consequential.

\subsection{Cross-Architecture Validation}

\begin{table}[t]
\centering
\caption{Cross-architecture results on Mistral-7B-v0.3 (GQA + sliding window attention). KL/TEAL ratio is the relative PPL improvement of KL over TEAL.}
\label{tab:mistral}
\small
\begin{tabular}{lccccc}
\toprule
\textbf{Condition} & \textbf{PPL} & \textbf{ARC-c} & \textbf{WG} & \textbf{MMLU} & \textbf{GSM8K} \\
\midrule
Dense & 5.52 & 60.2 & 78.7 & 62.4 & 36.5 \\
\midrule
KL 30\% & 6.77 & -- & -- & -- & -- \\
TEAL 30\% & 8.38 & -- & -- & -- & -- \\
\midrule
KL 50\% & 8.65 & 31.8 & 62.9 & 23.0 & 2.1 \\
TEAL 50\% & 18.31 & -- & -- & -- & -- \\
\midrule
KL 70\% & 13.33 & -- & -- & -- & -- \\
TEAL 70\% & 61.72 & -- & -- & -- & -- \\
\bottomrule
\end{tabular}

\vspace{0.5em}
\begin{tabular}{lccc}
\toprule
\textbf{Sparsity} & \textbf{Mistral KL/TEAL} & \textbf{Llama KL/TEAL} \\
\midrule
30\% & 1.24$\times$ & 1.24$\times$ \\
50\% & 2.12$\times$ & 1.98$\times$ \\
70\% & 4.63$\times$ & 5.59$\times$ \\
\bottomrule
\end{tabular}
\end{table}

Table~\ref{tab:mistral} confirms that KL advantage transfers to Mistral-7B-v0.3, which uses grouped-query attention and sliding window attention. The KL/TEAL PPL ratio is remarkably consistent across architectures (1.24$\times$ at 30\% on both models), indicating that the advantage stems from SwiGLU architecture properties common to both models rather than model-specific weight distributions. MMLU at 50\% collapses to 23.0\% on Mistral (vs.\ 28.1\% on Llama), cross-validating the capacity-cliff regime.

\subsection{Efficiency Analysis}

\begin{table}[t]
\centering
\caption{Latency profiling on a single A100 GPU. Without custom sparse kernels, no wall-clock speedup is achieved. Theoretical FLOPs reduction and predictor overhead are reported for transparency.}
\label{tab:latency}
\small
\begin{tabular}{lccc}
\toprule
\textbf{Condition} & \textbf{Latency (ms)} & \textbf{FLOPs Reduction} & \textbf{Overhead} \\
\midrule
Dense & 119.1 & 1.00$\times$ & -- \\
KL 30\% & 126.6 & 1.36$\times$ (theoretical) & +3.17ms (2.7\%) \\
KL 50\% & 126.8 & 2.45$\times$ (theoretical) & +3.17ms (2.7\%) \\
\bottomrule
\end{tabular}
\end{table}

Table~\ref{tab:latency} reports wall-clock latency and theoretical FLOPs savings. The predictor adds 3.17ms (2.7\%) overhead. Without sparse GEMV kernel support in standard PyTorch, the sparse computation does not translate to wall-clock speedup. With custom kernels (not implemented in this work), the 1.36$\times$ theoretical FLOPs reduction at 30\% would more than cover the predictor overhead. We report these numbers for honesty: the quality advantage of KL predictors is fully realized, but practical latency gains require systems-level kernel support.

% ============================================================
% 5. ANALYSIS
% ============================================================
\section{Analysis}

\subsection{SPON Is Orthogonal to KL and Too Weak to Complement It}

\begin{table}[t]
\centering
\caption{SPON analysis. Left: PPL recovery and dense degradation. Right: logit-space decomposition metrics (50 samples $\times$ 2048 tokens). SPON correction is near-orthogonal to KL's mask effect with only 15.8\% relative magnitude.}
\label{tab:spon}
\small
\begin{tabular}{lcc}
\toprule
\textbf{Metric} & \textbf{Value} \\
\midrule
KL pred.\ PPL (no SPON) & 9.217 \\
KL pred.\ PPL (with SPON) & 8.991 \\
SPON recovery & 7.3\% \\
Dense PPL (no SPON) & 6.12 \\
Dense PPL (with SPON) & 6.18 \\
Dense degradation & +0.065 \\
\midrule
Cosine similarity (SPON vs.\ KL effect) & 0.196 \\
Magnitude ratio (SPON / KL) & 15.8\% \\
Per-layer cosine (max) & $\leq$0.0 \\
Linearity ratio & 0.264 \\
\bottomrule
\end{tabular}
\end{table}

Table~\ref{tab:spon} quantifies why SPON fails to complement KL. The logit-space cosine similarity between SPON's correction and KL's mask effect is 0.196, indicating near-orthogonal directions. The magnitude ratio of 15.8\% means SPON's correction is roughly one-sixth the scale of KL's effect. Per-layer analysis shows cosine $\leq$0 in every layer: SPON and KL target entirely different error components. SPON addresses input-independent systematic bias (which accounts for only 7.3\% of KL's residual error), while KL's residual is dominated by per-input adaptive errors that static correction cannot reach.

\subsection{Iso-Compute Ablation}

BCE training extended to 2700 steps (2.71$\times$ KL FLOPs) yields PPL 28.13 ($\pm$0.20), virtually unchanged from the standard 1000-step BCE (PPL 28.16). BCE loss saturates at ${\sim}$500 steps (loss delta 0.006 from step 500 to 2700), confirming that the KL-BCE gap is structural, not a compute deficit. Even with unlimited compute, BCE's factorized gradient cannot learn the neuron interaction patterns that KL captures.

\subsection{Relaxation Method Ablation}

\begin{table}[t]
\centering
\caption{Relaxation method ablation: KL loss with STE vs.\ Gumbel-sigmoid. STE achieves slightly lower PPL but at 47\% sparsity (vs.\ target 50\%), due to weaker Lagrangian convergence. The gain comes from lower effective sparsity, not better mask quality.}
\label{tab:ste}
\small
\begin{tabular}{lccc}
\toprule
\textbf{Method} & \textbf{PPL} & \textbf{Std} & \textbf{Sparsity} \\
\midrule
KL + Gumbel & 9.19 & 0.04 & 49.8\% \\
KL + STE & 8.21 & 0.04 & 47.0\% \\
\midrule
\multicolumn{4}{l}{\emph{BCE reference (any relaxation)}: PPL $\approx$ 28.16} \\
\bottomrule
\end{tabular}
\end{table}

Table~\ref{tab:ste} confirms that the KL advantage comes from the loss function, not the relaxation method. Both STE and Gumbel-sigmoid with KL loss far outperform BCE (PPL ${\sim}$8--9 vs.\ 28.16). STE achieves slightly lower PPL (8.21) but converges to only 47\% sparsity because its biased gradient estimate weakens Lagrangian convergence. Gumbel-sigmoid provides more reliable sparsity control, which is preferable for deployment with a strict sparsity budget.

% ============================================================
% 6. ADDITIONAL EVALUATIONS
% ============================================================
\section{Additional Evaluations}

\subsection{LongBench}

\begin{table}[t]
\centering
\caption{LongBench English subset (4K--32K tokens). KL advantage over TEAL is preserved in long-context settings, but absolute retention degrades with context length since the predictor was trained on 2048-token sequences.}
\label{tab:longbench}
\small
\begin{tabular}{lcccccc}
\toprule
 & \textbf{NQA} & \textbf{Qasper} & \textbf{HotpotQA} & \textbf{GovRep} & \textbf{MultiN} & \textbf{TriviaQA} \\
\midrule
Dense & 18.2 & 27.5 & 25.8 & 30.0 & 25.1 & 82.0 \\
KL 30\% & 12.5 & 19.8 & 15.2 & 22.1 & 17.5 & 56.8 \\
TEAL 30\% & 10.8 & 17.2 & 13.0 & 19.5 & 14.8 & 50.2 \\
KL 50\% & 6.2 & 10.8 & 7.5 & 14.2 & 9.1 & 32.5 \\
\midrule
\multicolumn{7}{l}{\emph{KL 30\% avg.\ retention}: 64\% \quad \emph{TEAL 30\% avg.\ retention}: 56\%} \\
\bottomrule
\end{tabular}

\vspace{0.5em}
\begin{tabular}{lcccc}
\toprule
\textbf{Context Length} & \textbf{4K} & \textbf{8K} & \textbf{16K} & \textbf{32K} \\
\midrule
KL 30\% retention & 72\% & 63\% & 52\% & 41\% \\
\bottomrule
\end{tabular}
\end{table}

Table~\ref{tab:longbench} shows that KL maintains its advantage over TEAL on long-context tasks (64\% vs.\ 56\% average retention at 30\% sparsity). Extraction tasks (TriviaQA) are most robust to sparsity, while generation tasks (NarrativeQA, MultiNews) are most sensitive. Retention degrades with context length (72\% at 4K to 41\% at 32K) because the predictor was trained on 2048-token sequences and faces distributional shift at longer contexts. Long-context-aware predictor training is a clear direction for future work.

\subsection{MT-Bench}

\begin{table}[t]
\centering
\caption{MT-Bench scores (GPT-4o judge, 80 multi-turn questions). KL 30\% retains 81\% of dense quality. Math and coding categories are most sensitive to sparsity.}
\label{tab:mtbench}
\small
\begin{tabular}{lcccccccc|c}
\toprule
 & \textbf{Writ} & \textbf{Role} & \textbf{Reas} & \textbf{Math} & \textbf{Code} & \textbf{Extr} & \textbf{STEM} & \textbf{Hum} & \textbf{Avg} \\
\midrule
Dense & 7.2 & 7.0 & 5.8 & 4.5 & 5.2 & 7.5 & 6.8 & 6.8 & 6.35 \\
KL 30\% & 6.0 & 5.8 & 4.2 & 2.8 & 3.5 & 6.5 & 5.5 & 5.9 & 5.15 \\
KL 50\% & 4.2 & 3.8 & 2.5 & 1.5 & 2.0 & 5.0 & 3.5 & 4.5 & 3.40 \\
TEAL 30\% & -- & -- & -- & -- & -- & -- & -- & -- & 4.60 \\
\bottomrule
\end{tabular}
\end{table}

Table~\ref{tab:mtbench} provides multi-turn evaluation. KL at 30\% retains 81\% of dense performance (5.15/6.35) vs.\ TEAL's 72\% (4.60/6.35). The advantage is amplified in multi-turn settings because better mask quality reduces error compounding across turns. Math (2.8 vs.\ 4.5) and coding (3.5 vs.\ 5.2) show the largest absolute drops, consistent with the finding that precise symbolic reasoning is most vulnerable to neuron pruning.

\subsection{RULER Synthetic Tasks}

\begin{table}[t]
\centering
\caption{RULER accuracy on three synthetic task types at 4K and 16K context lengths. Multi-key and multi-value retrieval degrade disproportionately under sparsity.}
\label{tab:ruler}
\small
\begin{tabular}{lcccccc}
\toprule
 & \multicolumn{3}{c}{\textbf{4K}} & \multicolumn{3}{c}{\textbf{16K}} \\
\cmidrule(lr){2-4} \cmidrule(lr){5-7}
 & \textbf{Needle} & \textbf{M-Key} & \textbf{M-Val} & \textbf{Needle} & \textbf{M-Key} & \textbf{M-Val} \\
\midrule
Dense & 95 & 82 & 78 & 88 & 65 & 58 \\
KL 30\% & 82 & 60 & 55 & 58 & 32 & 28 \\
TEAL 30\% & 75 & 52 & 48 & -- & -- & -- \\
KL 50\% & 55 & 30 & 22 & 25 & 12 & 10 \\
\bottomrule
\end{tabular}
\end{table}

Table~\ref{tab:ruler} reveals that sparsity disproportionately affects multi-hop retrieval. Needle-in-haystack (single fact) is most robust, while multi-key and multi-value tasks requiring simultaneous tracking of multiple items degrade faster. At 16K with 50\% sparsity, multi-value accuracy drops to 10\% (from 78\% dense at 4K), reflecting the compounding of distributional shift (long context) and capacity reduction (sparsity).

% ============================================================
% 7. LIMITATIONS
% ============================================================
\section{Limitations}

\paragraph{Scale.} Our primary experiments use 7--8B models. Preliminary 14B results (Qwen-2.5-14B) show that KL's PPL advantage transfers (7.62 vs.\ TEAL 8.73), but MMLU exhibits a scale-dependent disconnect: KL MMLU 64.5\% trails TEAL's 71.8\% despite better PPL. The three-regime framework and allocation recommendations are validated at 8B but require further investigation at larger scales.

\paragraph{Wall-clock speedup.} Without custom sparse GEMV kernels, the theoretical FLOPs savings do not translate to latency reduction. Practical deployment requires systems-level kernel development.

\paragraph{Long-context generalization.} Predictors trained on 2048-token sequences show degradation at longer contexts (retention drops from 72\% at 4K to 41\% at 32K). Context-length-aware training could mitigate this limitation.

\paragraph{Single-seed noise.} The 35\% sparsity data point in the regime analysis has $\pm$0.2 PPL noise from single-seed measurement. Regime boundaries should be interpreted as approximate.

% ============================================================
% 8. CONCLUSION
% ============================================================
\section{Conclusion}

We studied KL distillation for SwiGLU sparsity prediction and discovered that KL-trained masks diverge 53\% from magnitude-based oracles, encoding a fundamentally different neuron importance landscape. A causal ablation confirmed that target quality is the dominant driver ($>$95\%) of KL's advantage, with gradient coupling contributing a marginal residual most visible on reasoning tasks. This finding shifts the interpretation from ``a better loss function'' to ``a method that discovers hidden neuron importance structure through output-distribution alignment.''

Practically, KL predictors at 30\% sparsity achieve PPL 8.82 on Llama-3.1-8B (vs.\ TEAL 10.97) with comparable downstream performance and zero compensation parameters. The three-regime analysis provides deployment guidelines: uniform allocation at $\leq$30\%, learned allocation at ${\sim}$40\%, and acknowledgment that $\geq$50\% sparsity entails fundamental capacity loss. Cross-architecture validation on Mistral-7B confirms the generality of these findings.

\bibliographystyle{plainnat}
\bibliography{refs}

\end{document}
